\subsection{$L^{2}$ parameter regularization}
\paragraph{Purpose}: This penalty is as well knows as \tB{weight decay}, it allows
to drives the weight closer to the origin.

\paragraph{Theory}
$\Omega(\bm{\theta}) = \dfrac{1}{2}\norm{\bm{w}}^{2}_{2}$\\
Then $\nabla_{\bm{w}}\tilde{J} = \alpha\bm{w} + \nabla_{\bm{w}}J$ and a gradient step of
size $\varepsilon$ becomes
\begin{center}
    $\bm{w} \leftarrow \left(1 - \varepsilon\alpha\right)\bm{w}
    - \varepsilon\nabla_{\bm{w}}J$
\end{center}
that is, the weights are \uB{shrunk by a constant factor} before each usual update.
Only the weights are penalized, the biases usually are not.\\
This is the MAP estimate with a Gaussian prior $\mathcal{N}\left(\bm{0},
\frac{1}{\alpha}\bm{I}\right)$ on the weights.

\subsection{$L^{1}$ parameter regularization}
\paragraph{Purpose}: it yields \tB{sparse} solutions, many weights being driven exactly to
zero.

\paragraph{Theory}
$\Omega(\bm{\theta}) = \norm{\bm{w}}_{1} = \su{i}{}\left|w_{i}\right|$, so that
$\nabla_{\bm{w}}\tilde{J} = \alpha\,\mathrm{sign}(\bm{w}) + \nabla_{\bm{w}}J$: the
penalty gradient has a \uB{constant magnitude}, whatever the size of the weight.\\
This is the MAP estimate with a Laplace prior on the weights.
